\subsection{Local intersection forms and the quartic tangent obstruction}

\begin{statement}{P4.51}{Proved}
Let $S$ be an irreducible reduced support in the torus
fibration of Part~4.2 with $b_1(S)>0$. Every local
Milnor fibre of a reduced local equation of $S$ has
identically zero rational intersection form. This applies
in particular to the remaining positive ruled support,
whose fundamental group is $\mathbb Z$.
\end{statement}
\begin{proof}
Let $U$ be a sufficiently small full tube and $F_t$ a
nearby compact smooth torus fibre. P4R.4 identifies the
homomorphism
\[
 \pi_1(F_t)=\mathbb Z^4\longrightarrow
 \pi_1(U)=\pi_1(S)
\]
and gives a finite cyclic cokernel, of order equal to
the fibre multiplicity. Its map on rational first
homology is therefore onto. Write
\[
 W=\ker\bigl(H_1(F_t;\mathbb Q)\longrightarrow
                  H_1(S;\mathbb Q)\bigr).
\]
Then $\dim W=4-b_1(S)\le3$.

Choose a coordinate ball about the point, with its
closed ball contained in $U$. The local full fibre
equation is $u g^m$, with $g$ reduced and $u$ a
holomorphic unit. We take a holomorphic $m$-th root of
$u$ on the ball and absorb it into $g$. The actual
nearby fibre in the ball is consequently the disjoint
union of the local Milnor fibres
$g=\zeta t^{1/m}$. Each such fibre $M$ is an oriented
codimension-zero submanifold with boundary of $F_t$.
Its inclusion into $U$ factors through the contractible
ball. Thus every loop of $M$ maps to zero in
$H_1(U;\mathbb Q)$, and the image subgroup
\[
 H=\operatorname{im}\bigl(\pi_1(M)\longrightarrow
                                 \mathbb Z^4\bigr)
\]
satisfies $H\otimes\mathbb Q\subset W$.

The covering criterion lifts $M\to F_t$ to the
connected cover $F_H\to F_t$ corresponding to $H$.
This cover is a $K(H,1)$, with $H$ free abelian. The
map on second rational homology is therefore exactly
\[
 \bigwedge^2(H\otimes\mathbb Q)\longrightarrow
              \bigwedge^2 H_1(F_t;\mathbb Q).
\]
Every class from $H_2(M;\mathbb Q)$ consequently maps
into $\bigwedge^2W$. The torus intersection pairing is
exterior multiplication followed by the orientation
isomorphism $\bigwedge^4H_1(F_t;\mathbb Q)=\mathbb Q$.
It vanishes on $\bigwedge^2W$ because $\dim W\le3$.

This ambient pairing is the local intersection
pairing on $M$. We represent finitely many cycles
inside its interior and make them transverse there;
the orientation-preserving inclusion into $F_t$
retains all intersection points and signs. Thus the
local pairing is zero. The symbolic exterior-square
matrix and its zero Gram matrix are recorded in
\texttt{code/part4/local\_intersection.py} and its
executed output.
\end{proof}

\begin{statement}{P4.52}{Proved}
In the remaining positive ruled support of Part~4.2,
the image of a contracted horizontal pole section has
ambient multiplicity exactly three. This restriction
imposes no generic-transverse-type hypothesis on
any component of the critical curve.
\end{statement}
\begin{proof}
We first give the comparison with a leading
homogeneous polynomial. Write the reduced equation
as $g=g_d+g_{d+1}+\cdots$, and suppose that $g_d$ is
reduced. Every finite collection of cycles in
$M_d=\{g_d=1\}$, together with its intersections,
transports into a local Milnor fibre of $g$.
Indeed, choose a compact neighbourhood of the cycle
representatives. Euler's identity shows that
$dg_d\ne0$ on $M_d$. The functions
\[
 G_\epsilon(x)=\epsilon^{-d}g(\epsilon x)
\]
converge with their derivatives to $g_d$ on that
neighbourhood. We solve
\[
 dG_\epsilon(V_\epsilon)
                     =-\partial_\epsilon G_\epsilon
\]
using the Hermitian gradient divided by its squared
norm. A cutoff supported in a slightly larger
compact neighbourhood, equal to one along the
trajectories under consideration, gives a flow for
small positive real $\epsilon$. It transports those
cycles to $G_\epsilon=1$. Holomorphic scaling by
$\epsilon$ then embeds them in $g=\epsilon^d$ inside
the fixed original ball, preserving orientations
and intersections. This compares smooth nonzero
level sets on compact neighbourhoods; no assertion
about equisingularity of the zero set is needed.
By P4.51 the pairing on $M_d$ is consequently zero.

For a reduced homogeneous polynomial its global
fibre and a sufficiently large compact Milnor fibre
have the same homology and scalar monodromy. One
can see this directly from the squared radius on
the affine fibre. This is a proper semialgebraic
function with finitely many critical values; above
the largest, its gradient supplies a product end.
A large invariant ball therefore gives a compact
core. Rescaling identifies that core with a small
local Milnor fibre. The Milnor-fibre CW dimension
bound gives zero homology above degree two, and
reducedness gives connectedness.

We need one exact variation identity. For such a
homogeneous polynomial, let $\mu_i$ be the
transverse plane-curve Milnor numbers at the
singular points of its reduced projective curve,
and put $M=\sum_i\mu_i$. Write $F$ for its compact
Milnor fibre and $\partial_2F$ for the boundary
part over the singular-link circles. Let
\[
 \begin{split}
 j&:H_2(F;\mathbb Q)\longrightarrow
                         H_2(F,\partial F;\mathbb Q),\\
 j_2&:H_2(F;\mathbb Q)\longrightarrow
                         H_2(F,\partial_2F;\mathbb Q).
 \end{split}
\]
Siersma's chain-level variation identity
\cite[\S3.6, pp.~451--452; \S7.1, p.~462]
{Siersma1991Variation} is
\[
                 j_2(T-I)=\operatorname{VAR}^{II}j.
                                                            \tag{4.52a}
\]
Here $T$ is the Milnor monodromy on $H_2(F;\mathbb Q)$.
The identity does not require the variation map to
be an isomorphism in dimension two. The dimension-two
Wang sequence is explicitly treated in that source,
\S6.3. Poincar\'e--Lefschetz duality identifies $j$
with the intersection form. Its vanishing makes
$\operatorname{im}(T-I)$ lie in $\ker j_2$.
The exact sequence of $(F,\partial_2F)$ identifies
that kernel with the image of $H_2(\partial_2F)$.
Each boundary piece is the mapping torus of a
transverse Milnor fibre. Its Wang sequence gives
\[
 H_2(\partial_2F;\mathbb Q)
       =\bigoplus_i\ker(A_i-I)\subset
                        \bigoplus_iH_1(F_i;\mathbb Q),
\]
where $A_i$ is the vertical monodromy. Thus zero
intersection form implies
\[
                 \operatorname{rank}(T-I)\le M.             \tag{4.52b}
\]
An empty $\partial_2F$ gives the same assertion with
$M=0$.

Suppose now that the reduced projective quartic
$D=\{g_4=0\}$ has a normalization component of
positive genus. By degree, either $D$ is an
irreducible quartic with normalization genus
$g\ge1$, or it is a smooth elliptic cubic and a
line. In the first case
$\sum\delta_i=3-g\le2$, so
$M=\sum(2\delta_i-r_i+1)\le4$.
In the second case B\'ezout gives total intersection
length three and at least one intersection point.
Hence $\sum\delta_i=3$ and
$\sum(r_i-1)\ge1$, giving $M\le5$.

The map $\{g_4=1\}\to\mathbb P^2\setminus D$ is
an unramified fourfold cover. The normalization
sequence and plane-curve genus formula give
\[
              \chi(F)=4\chi(\mathbb P^2\setminus D)
                      =4(7-M).
\]
The three nonidentity scalar monodromies act freely
on the compact invariant core, so their Lefschetz
numbers are zero. It follows by the character
formula for the cyclic group of order four that
\[
 H_0(F;\mathbb C)-H_1(F;\mathbb C)+H_2(F;\mathbb C)
       =(7-M)\,\mathbb C[\mathbb Z/4]
\]
as virtual representations. Every nontrivial
character therefore has multiplicity at least
$7-M$ in $H_2$. The finite-order monodromy is
semisimple, and hence
\[
             \operatorname{rank}(T-I)
                   \ge3(7-M)\ge6>M.                       \tag{4.52c}
\]
This contradicts (4.52b).

We apply this obstruction through the actual pole
comparison. Blow up the ambient origin, normalize
the strict transform, and take its minimal
resolution $T_1$. As in P4.39, this resolution maps
to the original normalization and dominates its
minimal resolution by point blowups. The strict
transform of the positive-genus horizontal section
$H$ retains pole coefficient two. If $E$ is the
ambient exceptional divisor, put
$n=\operatorname{ord}_H(E)>0$, and let $a\ge0$
be the coefficient of $H$ in the total conductor
divisor relative to the strict-transform dualizing
sheaf. Adjunction gives exactly
\[
                          2=(d-2)n+a.                     \tag{4.52d}
\]
Effectiveness follows from the exceptional
intersection matrix argument of P4R.6. An
exceptional positive-genus curve over the new
normalization has strictly positive coefficient,
by its adjunction equation. Thus $d\ge5$ is
impossible. If $d=4$, then $n=1,a=0$; the curve
$H$ is neither exceptional over the new
normalization nor over its conductor. The
normalization is an isomorphism at its generic
point. Its image is a multiplicity-one component
of the quartic exceptional curve, birational to
$H$. It has positive genus. Such a component has
degree three or four, leaving at most one linear
factor; the entire quartic polynomial is reduced.
The preceding contradiction excludes this case.

For clarity, an isolated critical point cannot
hide a contracted $H$. A generic small linear
perturbation of an isolated hypersurface critical
point has a positive finite number of Morse
critical points, the local degree of its gradient.
A small constant shift makes a chosen critical
value nonzero. On a fixed sufficiently small ball
the original zero-level boundary is smooth and
transverse; compactness retains the boundary
submersion over a uniform punctured value disc.
The finite-critical-value and path-transport
argument of P4.20 then carries the Morse
vanishing sphere, of square $-2$, into the
original local Milnor fibre. P4.51 excludes it.
This argument does not assume smoothness of the
normalization.

At a nonisolated point with $d=2$, the Hessian is
nonzero. P4R.15 gives the exact product equation
$X^2+Y^m$. For even $m$ its two analytic branches
are smooth. For odd $m$ its normalization is the
explicit smooth parametrization
$(Y,X,Z)=(u^2,\mathrm i u^m,Z)$.
Neither case has an exceptional curve in its
minimal normalization resolution. A point with
$d=1$ is smooth as well. Therefore the image of
$H$ has multiplicity exactly three.

The cyclic character calculation and all integer
bounds are independently recomputed in
\texttt{code/part4/quartic\_tangent.py}; no numerical
approximation enters the argument.
\end{proof}
